\DocumentMetadata{
  pdfversion=2.0,
  pdfstandard=ua-2,
  lang=en-US,
  tagging=on,
  tagging-setup={math/setup=mathml-SE,tabsorder=structure}
}
\documentclass[11pt,letterpaper]{article}
\usepackage[margin=0.68in,includefoot]{geometry}
\usepackage{newcomputermodern}
\usepackage{amsmath,amscd,mathtools}
\usepackage{tikz,adjustbox}
\providecommand{\square}{\mdlgwhtsquare}
\usepackage[hidelinks]{hyperref}
\hypersetup{
  pdftitle={Analysis PhD Exam, June 1996},
  pdfauthor={University of Florida Department of Mathematics},
  pdfsubject={Graduate mathematics examination},
  pdfdisplaydoctitle=true,
  bookmarksopen=true
}
\setcounter{secnumdepth}{0}
\setlength{\parindent}{0pt}
\setlength{\parskip}{0.28em}
\setlength{\emergencystretch}{3em}
\setlength{\leftmargini}{2.15em}
\setlength{\leftmarginii}{2.65em}
\setlength{\leftmarginiii}{3em}
\renewcommand{\labelenumi}{\textbf{\theenumi.}}
\pagestyle{empty}
\raggedbottom
\allowdisplaybreaks
\newenvironment{examproblems}[1]{%
  \begin{enumerate}%
  \setcounter{enumi}{#1}%
  \setlength{\topsep}{0.35em}%
  \setlength{\partopsep}{0pt}%
  \setlength{\itemsep}{0.48em}%
  \setlength{\parsep}{0.18em}%
}{\end{enumerate}}
\newenvironment{examsubparts}{%
  \begin{enumerate}%
  \setlength{\topsep}{0.18em}%
  \setlength{\partopsep}{0pt}%
  \setlength{\itemsep}{0.16em}%
  \setlength{\parsep}{0.1em}%
}{\end{enumerate}}
\newenvironment{examinstructions}{%
  \par\begingroup\itshape
}{%
  \par\endgroup\vspace{0.18em}
}
\makeatletter
\renewcommand\section{\@startsection{section}{1}{0pt}%
  {-1.25ex plus -.25ex minus -.1ex}%
  {0.55ex plus .1ex}%
  {\normalfont\large\bfseries}}
\renewcommand{\@maketitle}{%
  \newpage\null\vskip -1em%
  {\footnotesize\bfseries
    \href{https://www.ufl.edu/}{University of Florida}, \href{https://math.ufl.edu/}{Department of Mathematics}\par}%
  \vskip -0.5em%
  \tagstructbegin{tag=Title}%
    {\LARGE\bfseries\href{https://gma.math.ufl.edu/exams/analysis-phd/}{Analysis PhD Exam}, June 1996\par}%
  \tagstructend%
  \vskip 0.25em%
  \tagmcbegin{artifact}%
    \hrule height 1pt%
  \tagmcend%
  \vskip 1em%
}
\makeatother
\title{Analysis PhD Exam, June 1996}
\author{}
\date{}
\begin{document}
\maketitle
\thispagestyle{empty}
\begin{examinstructions}
Remark. In the sequel the word measure means positive, countably additive measure, unless otherwise stated.
\end{examinstructions}
\begin{examinstructions}
Instructions. Do all of the following problems.
\end{examinstructions}
\begin{examproblems}{0}
\item
Let \(f_n:\mathbb{R}\to\mathbb{R}\) be a sequence of Lebesgue-measurable functions which converges pointwise to a function~\(g\). Prove, directly from the definition of measurability, that~\(g\) is measurable.
\item
For \(f\in L^1(\mathbb{R})\), let \(f_t\) denote the translation of~\(f\) by~\(t\), i.e. \(f_t(x):=f(x-t)\). Show that the map \(t\mapsto f_t\) is a continuous map from~\(\mathbb{R}\) to \(L^1(\mathbb{R})\).
\item
Let~\(\mu\) denote the Lebesgue measure on~\(\mathbb{R}\). Let \(T:\mathbb{R}\to\mathbb{R}\) be defined by \(T(x)=x^3-x\). Define a Borel measure~\(\nu\) setting \(\nu(A):=\mu(T^{-1}A)\). Compute the Radon–Nikodym derivative \(\frac{d\nu}{d\mu}\).
\item
Let \(K_n\subset[0,1]\) be a Cantor set of Lebesgue measure larger than \(1-\frac{1}{n}\). Let \(V:=\bigcup_{n=1}^{\infty}K_n\). Prove that, for any null set~\(N\), the set \(V\setminus N\) is not a \(G_\delta\) set. [Hint: Baire Category Theorem]
\item
Let~\(\mu\) and~\(\nu\) be~\(\sigma\)-finite measures on the measurable space \((X,\mathcal{B})\). Prove, from first principles, that~\(\nu\) can be written as \(\nu=\nu_{\mathrm{ac}}+\nu_{\mathrm{sing}}\), with \(\nu_{\mathrm{ac}}\ll\mu\) and \(\nu_{\mathrm{sing}}\perp\mu\).
\item
Let \(x_n\) be a sequence in a Banach space~\(B\). Assume that the sequence \(x_n\) converges weakly to~\(x\).
\begin{examsubparts}
\item[\textnormal{(a)}]
Show that~\(x\) is the only weak limit of the sequence \(x_n\).
\item[\textnormal{(b)}]
Show that \(\sup_n\lVert x_n\rVert<\infty\). [Hint for (b): Consider the \(x_n\) as elements of the double dual \(B^{**}\).]
\end{examsubparts}
\item
Let \(K:[0,2\pi]\to\mathbb{R}\) be defined by \(K(x)=\frac14\sin(x)\). Show that for every \(h\in L^p([0,2\pi])\), there exists a unique solution~\(f\) to the equation 
\[
f+f*K=h.
\]
 If~\(h\) is \(C^\infty\), is it true that the solution~\(f\) is \(C^\infty\)? Justify your answer.
\item
A distribution~\(T\) is called harmonic if \(\Delta T=0\), where \(\Delta=\sum_{i=1}^n\frac{\partial^2}{\partial x_i^2}\). Show that if~\(T\) is a harmonic tempered distribution, then~\(T\) is a polynomial. [Hint: a distribution with support \(\{0\}\) is a finite sum of multiples of Dirac's~\(\delta\) and its derivatives.]
\end{examproblems}
\end{document}
